%======================================================================
\section{结论}
\label{sec:conclusion}
%======================================================================

我们将一个博弈论多智能体共识框架——AEIC 分解、三层 Stackelberg 机制与自适应协议路由——置于自动化学术评审上的"诊断后干预"验证之中。诊断结论是明确的：LLM 评审团存在羊群效应。独立评审携带系统性拒稿偏置（明确接收层论文接收率仅 14.4\%）；非结构化讨论触发负向棘轮，以 96:0 的单向翻转消灭 78\% 的少数派立场（$p{=}2.5{\times}10^{-29}$），评分方差坍缩 51\%（$p{=}3.9{\times}10^{-32}$），且从未改善决策一致率。羊群效应在三个模型家族中普遍存在，但其方向依赖于模型——DeepSeek-V3 与 R1 向拒稿棘轮化，GLM-4-Flash 向接收棘轮化——且最强的推理者遭受最大的多样性坍缩：能力不能抵御羊群效应，它只是选择了吸引子。

干预将机制类别映射到病理。聚合时校准——由 Stackelberg 优化自主发现为 $+1.04$ 偏置项——修复了拒稿偏置（明确接收层一致率 2\%$\to$38\%），而重设阈值的均值与之相当，因此结论是\emph{校准优先，形式不拘}。提示层面的行为调控无法打破棘轮（消灭率 78\%$\to$71\%），但保持了决策质量：话术敌不过从众先验，有效调控必须作用于机制或结构。逐论文路由以低 56\% 的成本匹配最优固定协议，将审议降级为按需资源。贯穿始终，行为健康指标被证明比决策准确率敏感数个数量级；我们倡议将少数派存活、方差保持与漂移作为一切多智能体判断系统的主要终点，将准确率降级为安全约束。

未来工作包括：在聚合时而非提示中保护少数派的机制级调控；混合模型家族的异质评审团（其相反的吸引子可能相互抵消）；更丰富的前瞻特征以弥合路由的 10 个百分点预言机差距；以及将审议健康指标迁移到生产级多智能体管道——推荐集成、多评审官评估、先审议后行动的智能体系统——作为可观测信号。

\begin{credits}
\subsubsection{\discintname}
The authors have no competing interests to declare.
\end{credits}
